\section{A group that needs two places}\label{sec:arithmetic}

\begin{proposition}\label{prop:twoplaces}
The group $\Gamma = \PGL_2(\Z[\frac12])$ is not discrete in $\PGL_2(\R)$ or in
$\PGL_2(\Q_2)$, but it is discrete in $\PGL_2(\R) \times \PGL_2(\Q_2)$.
\end{proposition}

\begin{proof}
The ring $\Z[\frac12]$ is dense in $\R$, so the upper triangular unipotent
matrices in $\Gamma$ are not discrete in $\PGL_2(\R)$. In $\PGL_2(\Q_2)$, the
infinite group $\PGL_2(\Z)$ lies in the compact group $\PGL_2(\Z_2)$, so
$\Gamma$ is not discrete there either. Finally, $\Z[\frac12]$ is discrete in
$\R \times \Q_2$, and so $\Gamma$ is discrete in the product; see
\cite[Chapter~II]{serre}, where $\Gamma$ acts on the product of the hyperbolic
plane and the $3$-regular tree.
\end{proof}

\begin{remark}
Proposition~\ref{prop:twoplaces} is why Sanov's group of
Theorem~\ref{thm:sanov} is \emph{geometric} while $\PGL_2(\Z[\frac12])$ is
\emph{arithmetic} in the sense of the Classification Conjecture.
\end{remark}
